% Gene-to-Phenotype Mapping Diagram
\begin{tikzpicture}[scale=0.8, every node/.style={scale=0.8}]

% Title
\node[text width=18cm, text centered] at (9,14) {
    \Large\textbf{DNA System: 16-Gene to 5-Phenotype Mapping in RCC}
};

% Gene categories
\node[draw, rounded corners, fill=dna!20, text width=16cm, minimum height=8cm] at (9,8) {};
\node[text width=16cm, text centered] at (9,12.5) {\textbf{16-Gene Mutation Profile}};

% Growth genes
\node[draw, rounded corners, fill=green!20, text width=3cm] at (2,11) {
    \textbf{Growth Genes}\\
    VHL\\
    PBRM1\\
    SETD2\\
    KDM5C\\
    PTEN
};

% Tumor suppressor
\node[draw, rounded corners, fill=red!20, text width=2.5cm] at (6,11) {
    \textbf{Suppressor}\\
    TP53
};

% Signaling
\node[draw, rounded corners, fill=blue!20, text width=2.5cm] at (10,11) {
    \textbf{Signaling}\\
    PIK3CA\\
    MET\\
    EGFR
};

% Immune evasion
\node[draw, rounded corners, fill=purple!20, text width=2.5cm] at (14,11) {
    \textbf{Evasion}\\
    CD274\\
    (PD-L1)
};

% Antigen presentation
\node[draw, rounded corners, fill=orange!20, text width=6cm] at (8,8.5) {
    \textbf{Antigen Presentation Complex}\\
    B2M, JAK1, JAK2, TAP1, TAP2, HLA-A\\
    (MHC Class I pathway)
};

% Mutation types legend
\node[draw, rounded corners, fill=yellow!10, text width=4cm] at (16,8.5) {
    \textbf{Mutation Types}\\
    0 = Wild Type\\
    1 = Loss of Function\\
    2 = Gain of Function\\
    3 = Amplification
};

% Mathematical transformation
\draw[very thick, ->] (9,6.5) -- (9,5.5);
\node[right] at (9.5,6) {\textbf{Mathematical Mapping}};

% Phenotype categories
\node[draw, rounded corners, fill=yellow!20, text width=16cm, minimum height=3cm] at (9,3) {};
\node[text width=16cm, text centered] at (9,4.5) {\textbf{5 Phenotypic Traits}};

% Proliferation chance
\node[draw, rounded corners, fill=green!30, text width=3cm, text centered] at (2,3) {
    \textbf{Proliferation}\\
    $P_{prolif} = f($VHL,$PBRM1,$SETD2,$KDM5C,$PTEN$)$
};

% Apoptosis resistance
\node[draw, rounded corners, fill=red!30, text width=2.5cm, text centered] at (5.5,3) {
    \textbf{Apoptosis}\\
    $P_{apopt} = f($TP53$)$
};

% Angiogenesis
\node[draw, rounded corners, fill=blue!30, text width=2.5cm, text centered] at (8.5,3) {
    \textbf{Angiogenesis}\\
    $A_{factor} = f($PIK3CA,$MET,$EGFR$)$
};

% Immune evasion
\node[draw, rounded corners, fill=purple!30, text width=2.5cm, text centered] at (11.5,3) {
    \textbf{PD-L1 Expr}\\
    $E_{PDL1} = f($CD274$)$
};

% Antigen presentation
\node[draw, rounded corners, fill=orange!30, text width=3cm, text centered] at (15,3) {
    \textbf{Ag Present}\\
    $P_{antigen} = f($MHC genes$)$
};

% Mapping arrows
\draw[thick, ->] (2,10) to[bend right=30] (2,3.8);
\draw[thick, ->] (6,10) to[bend right=20] (5.5,3.8);
\draw[thick, ->] (10,10) to[bend left=20] (8.5,3.8);
\draw[thick, ->] (14,10) to[bend left=30] (11.5,3.8);
\draw[thick, ->] (8,7.5) -- (15,3.8);

% Detailed mathematical formulations
\node[draw, rounded corners, fill=gray!10, text width=7cm] at (4,1) {
    \textbf{Example: Proliferation Calculation}\\
    For each growth gene $g$:\\
    $contrib_g = \begin{cases} 
    0.9 & \text{if } g = 1 \text{ (LoF)} \\
    1.1 & \text{if } g = 2 \text{ (GoF)} \\
    1.0 & \text{otherwise}
    \end{cases}$\\[0.2cm]
    $P_{proliferation} = \prod_{g \in \{VHL, PBRM1, ...\}} contrib_g$
};

\node[draw, rounded corners, fill=gray!10, text width=7cm] at (14,1) {
    \textbf{Example: Antigen Presentation}\\
    $damaged = \sum_{g \in MHC\_genes} \mathbb{1}[g \neq 0]$\\[0.1cm]
    $P_{antigen} = \max(0.1, 1.0 - 0.15 \times damaged)$\\[0.2cm]
    If all 6 MHC genes mutated:\\
    $P_{antigen} = 0.1$ (minimal presentation)
};

% Mutation dynamics
\node[draw, rounded corners, fill=blue!10, text width=8cm] at (9,0) {
    \textbf{Dynamic Mutation Process}\\
    Initial: 3 random mutations injected\\
    Duplication: daughter inherits parent mutations\\
    Additional: new mutations during simulation (rare)\\
    Clonal evolution: mutation accumulation over time
};

% Agent behavior influence arrows
\draw[very thick, ->, dashed] (2,2.2) to[bend right=40] (-1,0) node[left] {\textbf{Influences}};
\draw[very thick, ->, dashed] (15,2.2) to[bend left=40] (19,0) node[right] {\textbf{Agent Behavior}};

% Specific examples
\node[draw, rounded corners=3pt, fill=yellow!70] at (1,6) {Example 1};
\node[below] at (1,5.5) {\tiny VHL=1, CD274=2};
\node[below] at (1,5.2) {\tiny Fast growth};
\node[below] at (1,4.9) {\tiny High PD-L1};

\node[draw, rounded corners=3pt, fill=yellow!70] at (17,6) {Example 2};
\node[below] at (17,5.5) {\tiny TP53=1, B2M=1};
\node[below] at (17,5.2) {\tiny Apoptosis resist};
\node[below] at (17,4.9) {\tiny Immune invisible};

% Hamming distance for TCR matching
\node[draw, rounded corners, fill=orange!40, text width=4cm] at (1,9) {
    \textbf{TCR Matching}\\
    Hamming distance\\
    between TCR and\\
    tumor neoantigen\\
    $d_H \leq 3$ = match
};

\end{tikzpicture}nd{tikzpicture}
